\clearpage
\VerdierPhysicalPageAnchor{0168}%
\begingroup\small\setlength{\parskip}{0pt}\setlength{\emergencystretch}{2.5em}
\VerdierSemanticAnchor{ch4-4-2-4-sequence}
\noindent\textit{Le composé des morphismes :}
\[
E\!\begin{bmatrix}n-1\\ f_5,f_6,f_7\end{bmatrix}
\xrightarrow{d_1}
E\!\begin{bmatrix}n\\ f_3,f_4,f_5\end{bmatrix}
\xrightarrow{d_2}
E\!\begin{bmatrix}n+1\\ f_1,f_2,f_3\end{bmatrix}
\]
\textit{est nul. On a un isomorphisme canonique :}
\VerdierSemanticAnchor{ch4-4-2-4-1}
\[
\operatorname{Ker}(d_2)\big/\operatorname{Im}(d_1)
\simeq
E\!\begin{bmatrix}n\\ f_3f_2,f_4,f_6f_5\end{bmatrix}.
\tag{4.2.4.1}
\]
Résulte immédiatement de (4.2.3.3) et (4.2.3.1).

\VerdierSemanticAnchor{ch4-4-2-5}
\medskip\noindent\textbf{4.2.5.}
Les morphismes $d^n(\omega_5)$,
$\omega_5\in\operatorname{Ob}\mathcal{F}\ell_5(\mathcal J)$, sont appelés
les \textit{différentielles} (relatives à l'objet spectral considéré). La
famille des foncteurs (4.2.3.1)
\[
\omega_3\longmapsto
E\!\begin{bmatrix}n\\ \omega_3\end{bmatrix},
\qquad
\omega_3\in\operatorname{Ob}\mathcal{F}\ell_3(\mathcal J),
\]
munis des différentielles (4.2.3.2) $d^n(\omega_5)$,
$\omega_5\in\operatorname{Ob}\mathcal{F}\ell_5(\mathcal J)$, est appelée
la \textit{famille spectrale} associée à l'objet spectral
$(H^n(f),\delta)$.

\VerdierSemanticAnchor{ch4-4-2-6}
\medskip\noindent\textbf{4.2.6.}
Soit
$\omega_4=(\mathord\bullet\xrightarrow{f_1}\mathord\bullet
\xrightarrow{f_2}\mathord\bullet\xrightarrow{f_3}\mathord\bullet
\xrightarrow{f_4}\mathord\bullet)$
un objet de $\mathcal{F}\ell_4(\mathcal J)$. On en déduit un diagramme dans
$\mathcal{F}\ell_3(\mathcal J)$ :
\VerdierSemanticAnchor{ch4-4-2-6-fl3-diagram}
\[
\begin{tikzcd}[column sep=2.6em,row sep=2.15em]
{}\arrow[r,"\mathrm{id}"]\arrow[d,"f_1"']&
{}\arrow[r,"\mathrm{id}"]\arrow[d,"f_1"']&
{}\arrow[r,"\mathrm{id}"]\arrow[d,"f_1"']&
{}\arrow[r,"f_1"]\arrow[d,"f_2f_1"']&
{}\arrow[d,"f_2"]\\
{}\arrow[r,"\mathrm{id}"]\arrow[d,"f_2"']&
{}\arrow[r,"\mathrm{id}"]\arrow[d,"f_2"']&
{}\arrow[r,"f_2"]\arrow[d,"f_3f_2"']&
{}\arrow[r,"\mathrm{id}"]\arrow[d,"f_3"']&
{}\arrow[d,"f_3"]\\
{}\arrow[r,"\mathrm{id}"]\arrow[d,"f_3"']&
{}\arrow[r,"f_3"]\arrow[d,"f_4f_3"']&
{}\arrow[r,"\mathrm{id}"]\arrow[d,"f_4"']&
{}\arrow[r,"\mathrm{id}"]\arrow[d,"f_4"']&
{}\arrow[d,"f_4"]\\
{}\arrow[r,"f_4"]&
{}\arrow[r,"\mathrm{id}"]&
{}\arrow[r,"\mathrm{id}"]&
{}\arrow[r,"\mathrm{id}"]&{}
\end{tikzcd}
\]
D'où, en appliquant le foncteur
$E\!\begin{bmatrix}n\\ \omega_3\end{bmatrix}$, trois suites exactes de
$\mathcal A$ :
\VerdierSemanticAnchor{ch4-4-2-6-exact-sequences}
\[
0\longrightarrow
B\!\begin{bmatrix}n\\ f_2,f_4f_3\end{bmatrix}
 \Big/
B\!\begin{bmatrix}n\\ f_2,f_3\end{bmatrix}
\longrightarrow
E\!\begin{bmatrix}n\\ f_1,f_2,f_3\end{bmatrix}
\longrightarrow
E\!\begin{bmatrix}n\\ f_1,f_2,f_4f_3\end{bmatrix}
\longrightarrow0,
\]
\[
0\longrightarrow
E\!\begin{bmatrix}n\\ f_1,f_2,f_4f_3\end{bmatrix}
\longrightarrow
E\!\begin{bmatrix}n\\ f_1,f_3f_2,f_4\end{bmatrix}
\longrightarrow
E\!\begin{bmatrix}n\\ f_2f_1,f_3,f_4\end{bmatrix}
\longrightarrow0,
\]
\endgroup
